\chapter{Spis skrótów i symboli}

\begin{itemize}
\item[SPL] poziom natężenia dźwięku (ang. \english{Sound Pressure Level})
\item[ANC] aktywna redukcja hałasu (ang. \english{Active Noise Control})
\item[ADC] przetwornik analogowo-cyfrowy (ang. \english{Analog To Digital
converter}) 
\item[DAC] przetwornik cyfrowo-analogowy (ang. \english{Digital To Analog
converter})
\item[SMPS] zasilacz impulsowy (ang. \english{Switching-mode Power Supply})
\item[PSU] sekcja zasilania (ang. \english{Power Supply Unit})
\item[AP] zasilanie sekcji analogowych (ang. \english{Analog Power}) 
\item[DP] zasilanie sekcji cyfrowych (ang. \english{Digital Power})
\item[REF] zasilanie referencyjne (ang. \english{Reference Voltage})
\item[eMMC] zewnętrzna pamięć nieulotna (ang. \english{External MultiMedia Card})
\item[EMIF] filtracja zakłóceń wysokiej częstotliwości (ang.
    \english{Electromagnetic Interference Filtering})
\item[DSP] przetwarzanie sygnałów cyfrowych lub procesor sygnałowy (ang.
    \english{Digital Signal Processing}, ang. \english{Digital Signal
    Processor})
\item[McASP] wielokanałowy szeregowy interfejs audio (ang.
\english{Multi-channel Audio Serial Port})
\item[EDA] oprogramowanie wspierające projektowanie układów elektroniczych (ang.
    \english{Electronic Design Automation})
\item[FIR] filtr o skończonej odpowiedzi impulsowej (ang. \english{Finite
Impulse Response})
\item[SIMD] architektura komputerowa obejmująca systemy, w których
przetwarzanych jest wiele strumieni danych w oparciu o pojdeynczy strumień
rozkazów (ang. \english{Single Instruction, Multiple Data})
\item[LMS] algorytm adaptacyjny oparty na gradientowej metodzie najmniejszego
    spadku (ang. \english{Least Mean Squares}) 
\item[FPU] jednostka zmiennoprzecinkowa (ang. \english{Floating Point Unit})
\item[ALSA] system obsługi dźwięku dla systemów GNU/Linux (ang.
\english{Advanced Linux Sound Architecture})
\item[FXLMS] algorytm adaptacyjny oparty na gradientowej metodzie najmniejszego
    spadku z filtracją wejścia (ang. \english{Filtered X Least Mean Squares}) 
\item[NLMS] algorytm adaptacyjny oparty na gradientowej metodzie najmniejszego
    spadku z normalizacją danych wejściowychj (ang. \english{Normalized Least
    Mean Squares})
\item[FXNLMS] algorytm adaptacyjny oparty na gradientowej metodzie najmniejszego
spadku z filtracją i normalizacją danych wejściowych (ang. \english{Filtered X
Normalized Least Mean Squares})
\item[LFSR] rejestr przesuwny z liniowym sprzężeniem zwrotnym (ang.
\english{Linear Feedback Shift Register})
\item[LCG] linowy generator kongruencjalny (ang. \english{Linear Congruential
Generator})
\item[PSD] gęstość widmowa mocy (ang. \english{Power Spectral Density})

 
\end{itemize}
